\scanpage{13}
\diagram{\begin{tikzcd}
\unsure{0}\arrow[r]&G^0\arrow[r,dash]\arrow[d,equal]&G\arrow[r]\arrow[d,equal]&\widehat{\mathbf Z}\arrow[r]\arrow[d,equal]&0\\
&D(\Gamma/\text{racines de }1)&D(\Gamma)\arrow[r,leftrightarrow]\arrow[d]&D(\text{racines de }1)&\\
&&\text{nbres de Weil}&&
\end{tikzcd}}
\margin{groupe de Galois mot. de $\mathbf F_p$}
\diagram{\begin{tikzcd}
U=\operatorname{Ker}(\epsilon|G^0)\arrow[dd,"\unsure{D}"]&\simeq\mu_2\times U^0\\
\deleted{D(\Gamma/\unsure{\epsilon^*|i^*}\mathbf Z)}&\deleted{\simeq D(\mathbf Z/2\mathbf Z)}\\
\Gamma/(\text{racines de }1).\unsure{\epsilon^*|i^*}(\mathbf Z)&\mathbf Z/2\mathbf Z\times\Gamma^0/\text{racines de }1\arrow[uu,leftrightarrow]
\end{tikzcd}}
\margin{$\Gamma^0=$ unités de Weil (nbres \unsure{algébriques} de val. absolue $1$ en \uline{toutes} les places $\ne p$).}

\diagram{\begin{tikzcd}
G_m\arrow[r,dash,"i"]&G\arrow[r,dash,"\epsilon"]&G_m\\
\mathbf Z&\Gamma\arrow[l,"i^*"']&\mathbf Z\arrow[l,"\epsilon^*"']\\
|\lambda|=p^{i^*(\lambda)/2}&&\epsilon^*(1)=p
\end{tikzcd}}
\[
\epsilon i(\lambda)=\lambda^2
\]

$\unsure{\mathfrak f|F}\in G(Q)$ canonique, défini par $\Gamma\subset\overline Q^*$ ($\overline Q$, n'importe quelle ext. algébrique de $Q$ par rapport à laquelle on a défini $\Gamma$)

\diagram{\begin{tikzcd}
&(G_{\mathbf F_p}^0)_{\mathrm{red}}\arrow[rr]&&(G^0)_{Q_p}\arrow[d,dash]&\\
G_{\mathbf F_p}\arrow[r]&(G_{\mathbf F_p})_{\mathrm{red}}\arrow[rr,"\alpha","(\text{épi})"']\arrow[dr]\arrow[d,dash]&&(G)_{Q_p}\arrow[dr]&\\
&\begin{array}{c}\text{\unsure{groupe} \unsure{associé aux}}\\F\text{-cristaux}\\\text{sur }\mathbf F_p\end{array}\arrow[d,equal]&\widehat{\mathbf Z}\arrow[rr,equal]&&\widehat{\mathbf Z}\\
&D(\overline Q_p^*)&&&
\end{tikzcd}}
\margin{enveloppe algébrique de $\mathbf Z$}

\diagram{\begin{tikzcd}
Q\arrow[d,equal]&\overline Q_p^*\arrow[l,"\begin{array}{c}\text{val. }p\text{-adique}\\\text{normalisée}\\\downarrow\end{array}"']\arrow[r,leftrightarrow,"\alpha^*","\text{l'inclusion canonique}"']\arrow[d]&\Gamma\arrow[d]\\
Q&\overline Q_p^*/\text{racines de }1\arrow[l]\arrow[r]&\Gamma/\text{racines de }1
\end{tikzcd}}

$M$ motif \inserted{simple} sur $\mathbf F_p$, $A_M=\operatorname{End}(M)$ \deleted{$\unsure{E}$} \uline{centre} $Z_M$, \uline{engendré par} $\unsure{\mathfrak f_M|F_M}\in A$
\[
A\otimes_Q Q_p\simeq\operatorname{End}(T_p(M))\quad(\text{par Tate-Honda si }M\text{ une VA})
\]
\deleted{\unsure{En}}
\diagram{\begin{tikzcd}
{}\\
|[draw,ellipse]|\inserted{A_M\text{ splitte sur }Q_\ell\text{ pour }\ell\ne p\ (\ell\text{ fini})}\arrow[u,bend left]
\end{tikzcd}}
\margin{Les invariants des composants simples de}
\mspage{12}
